\documentclass[14pt,a4paper,oneside]{extbook}
\usepackage[utf8]{vietnam}
\usepackage{amsmath,amssymb,mathrsfs,fancyhdr,enumerate,multirow,makecell}
\usepackage{tkz-euclide,tikz-3dplot,tikz,tkz-tab}
\usetikzlibrary{arrows.meta,calc,intersections,angles,quotes,backgrounds,shapes.geometric,patterns,through,shadows,positioning}
\usepackage{yhmath}
\usepackage[top=0.7cm, bottom=0.7cm, left=1.0cm, right=1.0cm, headheight=16pt, headsep=3mm, footskip=5mm]{geometry}
\usepackage[hidelinks,unicode]{hyperref}
\usepackage{qrcode}
\usepackage[many]{tcolorbox}
\usepackage{fontawesome}

\pagestyle{fancy}
\fancyhf{}
\fancyhead[L]{\footnotesize\bfseries\color{teal!85!black}\faBook~SỔ TAY KIẾN THỨC TOÁN 5}
\fancyhead[R]{\footnotesize\bfseries\color{blue!85!black}Lớp Toán Thầy Lĩnh -- loptoanthaylinh.vn}
\fancyfoot[C]{\footnotesize\bfseries\color{gray!80!black}-- Sổ tay kiến thức: Bài 69. Ôn tập các phép tính --}
\renewcommand{\headrulewidth}{0.4pt}

\begin{document}

% --- BANNER TIÊU ĐỀ BÀI HỌC ---
\noindent
\begin{tcolorbox}[
  enhanced,
  colback=teal!8!white,
  colframe=teal!80!black,
  arc=3mm,
  boxrule=1.1pt,
  width=\linewidth,
  halign=center,
  top=2.2mm, bottom=2.2mm,
  left=3mm, right=3mm,
  drop fuzzy shadow
]
  {\large\bfseries\color{teal!95!black}\faBookmark\ BÀI 69. ÔN TẬP CÁC PHÉP TÍNH (CỘNG, TRỪ, NHÂN, CHIA)}\\[2pt]
  {\scriptsize\color{black!75}\itshape Quy tắc tính toán, tính chất giao hoán, kết hợp, phân phối và thứ tự biểu thức trên 1 trang A4}
\end{tcolorbox}

\vspace{3.2mm}
% --- KHỐI 1: CÁC PHÉP TÍNH VỚI PHÂN SỐ VÀ STP ---
\noindent
\begin{minipage}[t]{0.54\textwidth}
  \begin{tcolorbox}[
    enhanced,
    title={\footnotesize\bfseries\color{white}\faCalculator\ 1. Quy tắc bốn phép tính},
    colback=blue!2!white,
    colframe=blue!75!black,
    arc=2.5mm, boxrule=0.8pt, top=1.8mm, bottom=1.8mm, left=2.2mm, right=2.2mm
  ]
    \footnotesize
    $\bullet$ \textbf{Cộng/Trừ STP:} Đặt tính dấu phẩy thẳng cột.\\
    $\bullet$ \textbf{Nhân STP:} Nhân như số tự nhiên, đếm tổng chữ số TP ở cả hai thừa số để đặt dấu phẩy.\\
    $\bullet$ \textbf{Chia STP:} Dời dấu phẩy ở số chia sang tận cùng bên phải, đồng thời dời phẩy ở số bị chia tương ứng.
  \end{tcolorbox}
\end{minipage}%
\hfill
\begin{minipage}[t]{0.44\textwidth}
  \begin{tcolorbox}[
    enhanced,
    title={\footnotesize\bfseries\color{white}\faBolt\ 2. Tính chất tính nhanh},
    colback=teal!2!white,
    colframe=teal!75!black,
    arc=2.5mm, boxrule=0.8pt, top=1.8mm, bottom=1.8mm, left=2mm, right=2mm
  ]
    \footnotesize
    $\bullet$ \textbf{Giao hoán:} $a + b = b + a$; $a \times b = b \times a$\\
    $\bullet$ \textbf{Kết hợp:} $(a+b)+c = a+(b+c)$\\
    $\bullet$ \textbf{Phân phối:}\\
    \centerline{$\mathbf{a \times (b + c) = a \times b + a \times c}$}\\
    Đặt thừa số chung để nhóm số tròn chục.
  \end{tcolorbox}
\end{minipage}

\vspace{3mm}

% --- KHỐI 3: SƠ ĐỒ TƯ DUY ---
\noindent
\begin{tcolorbox}[
  enhanced,
  title={\footnotesize\bfseries\color{white}\faSitemap\ Sơ đồ tư duy: Bốn phép tính và Thứ tự thực hiện},
  colback=yellow!2!white,
  colframe=teal!85!black,
  arc=3mm, boxrule=0.9pt, top=1.8mm, bottom=1.8mm, left=1.5mm, right=1.5mm,
  halign=center
]
  \begin{tikzpicture}[
    scale=0.92,
    every node/.style={transform shape},
    root/.style={rectangle, rounded corners=5pt, fill=teal!85!black, text=white, font=\bfseries\footnotesize, inner sep=5pt, align=center},
    b1/.style={rectangle, rounded corners=4pt, fill=blue!12, draw=blue!75!black, text=blue!90!black, font=\bfseries\scriptsize, inner sep=4pt, align=center},
    b2/.style={rectangle, rounded corners=4pt, fill=purple!12, draw=purple!75!black, text=purple!90!black, font=\bfseries\scriptsize, inner sep=4pt, align=center},
    b3/.style={rectangle, rounded corners=4pt, fill=orange!12, draw=orange!75!black, text=orange!90!black, font=\bfseries\scriptsize, inner sep=4pt, align=center},
    b4/.style={rectangle, rounded corners=4pt, fill=green!15, draw=green!65!black, text=green!80!black, font=\bfseries\scriptsize, inner sep=4pt, align=center}
  ]
    \node[root] (core) at (0,0) {\faCalculator\ BỐN PHÉP TÍNH};
    \node[b1] (n1) at (-5.3, 0.75) {\textbf{CỘNG / TRỪ:}\\[1.5pt]Thẳng hàng dấu phẩy / Quy đồng mẫu};
    \node[b2] (n2) at (5.3, 0.75) {\textbf{NHÂN:}\\[1.5pt]Tử $\times$ Tử, Mẫu $\times$ Mẫu / Đếm chữ số TP};
    \node[b3] (n3) at (-5.3, -0.75) {\textbf{CHIA:}\\[1.5pt]Nhân với phân số đảo / Bỏ phẩy ở số chia};
    \node[b4] (n4) at (5.3, -0.75) {\textbf{THỨ TỰ BIỂU THỨC:}\\[1.5pt]Ngoặc trước $\rightarrow$ Nhân/Chia $\rightarrow$ Cộng/Trừ};

    \draw[teal!85!black, thick, -{Stealth[length=2mm]}] (core.west) -| (n1.south);
    \draw[teal!85!black, thick, -{Stealth[length=2mm]}] (core.east) -| (n2.south);
    \draw[teal!85!black, thick, -{Stealth[length=2mm]}] (core.west) -| (n3.north);
    \draw[teal!85!black, thick, -{Stealth[length=2mm]}] (core.east) -| (n4.north);
  \end{tikzpicture}
\end{tcolorbox}

\vspace{3mm}

% --- KHỐI 4: CẶP SỐ VÀNG & BẪY TRẮC NGHIỆM ---
\noindent
\begin{minipage}[t]{0.54\textwidth}
  \begin{tcolorbox}[
    enhanced,
    title={\footnotesize\bfseries\color{white}\faKey\ Các cặp số vàng nhân nhẩm siêu tốc},
    colback=green!3!white,
    colframe=green!60!black,
    arc=2.5mm, boxrule=0.8pt, top=1.8mm, bottom=1.8mm, left=2.2mm, right=2.2mm
  ]
    \footnotesize
    $\bullet$ $2 \times 5 = 10 \quad|\quad 4 \times 25 = 100 \quad|\quad 8 \times 125 = 1\,000$\\
    $\bullet$ $0,5 = \frac{1}{2} \quad|\quad 0,25 = \frac{1}{4} \quad|\quad 0,125 = \frac{1}{8}$\\
    $\bullet$ Chia cho $0,5$ tức là nhân với $2$; chia cho $0,25$ là nhân với $4$.
  \end{tcolorbox}
\end{minipage}%
\hfill
\begin{minipage}[t]{0.44\textwidth}
  \begin{tcolorbox}[
    enhanced,
    title={\footnotesize\bfseries\color{white}\faExclamationTriangle\ Mẹo tránh bẫy trắc nghiệm},
    colback=red!2!white,
    colframe=red!70!black,
    arc=2.5mm, boxrule=0.8pt, top=1.8mm, bottom=1.8mm, left=2.2mm, right=2.2mm
  ]
    \footnotesize
    $\bullet$ \textbf{Bẫy chia cho 0:} Phép chia cho số $0$ hoàn toàn vô nghĩa!\\
    $\bullet$ Trong biểu thức chỉ có phép nhân và chia: Thực hiện lần lượt từ \textbf{trái sang phải}, không được tiện tay làm nhân trước!
  \end{tcolorbox}
\end{minipage}

\end{document}
